\section{人类位置：从效率焦虑走向 agency 设计}

\subsection{教育、职业与“能力断层”风险}

访谈里一个高频问题是：初级工程师如何在 AI 高自动化环境下成长为专家。Sebastian 的回答很朴素但重要：保留“离线深度学习时间”。如果所有困难都外包给模型，人会失去构建抽象和调试直觉的机会。

\begin{importantbox}{AI 时代的人才培养建议（可执行版）}
\begin{itemize}[leftmargin=1.5em]
    \item 每天固定无 AI 时段，练习从零推导与实现；
    \item 把 AI 输出当草案，不当答案，强制做最小复现实验；
    \item 训练“规格表达 + 结果审查 + 失败复盘”三项能力。
\end{itemize}
\end{importantbox}

\subsection{意义、社区与长期社会结构}

后段讨论超出技术本身：即便 AI 带来更高生产率，人类依旧需要 agency、社区与被需要感。Lex 和两位嘉宾都强调，宏观效率提升不能自动解决个体失业与身份焦虑。未来高价值内容将更偏向“有真人痕迹”的创作和线下体验。

\begin{knowledgebox}{“human premium”可能上升的三个领域}
\begin{itemize}[leftmargin=1.5em]
    \item 线下协作与现场服务：不可复制的在场性；
    \item 高信任内容：可追溯作者身份与责任链；
    \item 长周期关系型工作：需要持续情感与组织承诺。
\end{itemize}
\end{knowledgebox}

\begin{warningbox}{技术讨论忽视个体体验会带来治理反噬}
若公共叙事只强调 GDP 与效率，而忽视个体失业、技能贬值与心理负担，社会对技术的反弹会加剧，最终反过来拖慢技术落地与制度建设。
\end{warningbox}

\subsection{本章小结}

AI 时代的关键不只是“能做什么”，而是“由谁决定做什么、承担什么后果”。agency 设计将成为技术与社会之间的接口能力。

